\section{Socket-Heating Protection Validation}
\label{sec:ch4-socket-heating-protection}

Socket-heating protection validation evaluated the relay-facing estimated
socket temperature. The important distinction is between four temperatures:
the direct NTC temperature, the expected normal socket temperature, the
estimated socket temperature used by protection, and the reference terminal
temperature. The trip threshold applies to the estimated socket temperature,
while the reference terminal temperature shows how hot the socket contact area
became during the same trial.
The plotted trends in this section were generated from the full protection
validation records in Appendix~\ref{app:d}.

The water-heater-only protection trend is shown in
Figure~\ref{fig:ch4-protect-water-heater}.
\begin{figure}[H]
  \centering
  \socketProtectionPlot{1}{240}
  \caption{Water heater temperature trend during socket-heating protection validation.}
  \label{fig:ch4-protect-water-heater}
\end{figure}

Figure~\ref{fig:ch4-protect-water-heater} shows that the low-current case did
not trip for either socket condition within the observation window. The
corroded terminal still ended hotter than the normal terminal, but the
estimated socket temperature stayed below the 40°C trip threshold.

The water heater plus one halogen lamp protection trend is shown in
Figure~\ref{fig:ch4-protect-one-lamp}.
\begin{figure}[H]
  \centering
  \socketProtectionPlot{2}{240}
  \caption{Water heater plus one halogen lamp trend during protection validation.}
  \label{fig:ch4-protect-one-lamp}
\end{figure}

Figure~\ref{fig:ch4-protect-one-lamp} shows that the normal socket did not
trip, while the corroded socket reached a protection-relevant state after a
longer heating period. This supports the estimator's role in distinguishing
slow thermal buildup from immediate current overload.

The water heater plus two halogen lamps protection trend is shown in
Figure~\ref{fig:ch4-protect-two-lamps}.
\begin{figure}[H]
  \centering
  \socketProtectionPlot{3}{120}
  \caption{Water heater plus two halogen lamps trend during protection validation.}
  \label{fig:ch4-protect-two-lamps}
\end{figure}

Figure~\ref{fig:ch4-protect-two-lamps} shows that the corroded socket tripped
much earlier than the normal socket at approximately the same current. The
reference terminal temperature was already above the estimated value at trip,
which confirms that the estimator acts as an early proxy for a hotter internal
condition.

The water heater plus three halogen lamps protection trend is shown in
Figure~\ref{fig:ch4-protect-three-lamps}.
\begin{figure}[H]
  \centering
  \socketProtectionPlot{4}{58}
  \caption{Water heater plus three halogen lamps trend during protection validation.}
  \label{fig:ch4-protect-three-lamps}
\end{figure}

Figure~\ref{fig:ch4-protect-three-lamps} shows fast thermal development in
the corroded case. The relay acted within the short test window, while the
reference terminal temperature reached 55.0°C.

The water heater plus four halogen lamps protection trend is shown in
Figure~\ref{fig:ch4-protect-four-lamps}.
\begin{figure}[H]
  \centering
  \socketProtectionPlot{5}{46}
  \caption{Water heater plus four halogen lamps trend during protection validation.}
  \label{fig:ch4-protect-four-lamps}
\end{figure}

Figure~\ref{fig:ch4-protect-four-lamps} shows the most severe heating case.
The corroded socket tripped after 5min when the estimated socket temperature
reached 40.0°C, while the reference terminal temperature had already
reached 56.0°C.

Table~\ref{tab:ch4-heating-protection} summarizes the protection validation
results.
\begin{table}[H]
  \centering
  \caption{Socket-heating protection validation results}
  \label{tab:ch4-heating-protection}
  \small
  \setlength{\tabcolsep}{3pt}
  \begin{tabular}{@{}p{0.20\textwidth}p{0.10\textwidth}p{0.07\textwidth}p{0.09\textwidth}p{0.08\textwidth}p{0.10\textwidth}p{0.10\textwidth}p{0.07\textwidth}@{}}
    \toprule
    Load case & Socket & Mean A & NTC 38.2 & Est. 40 & Est. temp & Terminal & Relay \\
    \midrule
    Water heater & Normal & 5.17 & -- & -- & 32.7 & 33.8 & No \\
    Water heater & Corroded & 4.85 & -- & -- & 36.7 & 38.5 & No \\
    Water heater + 1 lamp & Normal & 7.11 & -- & -- & 38.1 & 40.3 & No \\
    Water heater + 1 lamp & Corroded & 7.10 & 100 & -- & 39.8 & 42.4 & Yes \\
    Water heater + 2 lamps & Normal & 9.06 & 100 & 120 & 40.0 & 43.2 & Yes \\
    Water heater + 2 lamps & Corroded & 9.04 & 48 & 52 & 40.0 & 44.1 & Yes \\
    Water heater + 3 lamps & Normal & 11.07 & 54 & -- & 39.9 & 43.9 & Yes \\
    Water heater + 3 lamps & Corroded & 11.14 & 7 & -- & 39.7 & 55.0 & Yes \\
    Water heater + 4 lamps & Normal & 13.02 & 44 & -- & 39.9 & 45.0 & Yes \\
    Water heater + 4 lamps & Corroded & 13.08 & 5 & 5 & 40.0 & 56.0 & Yes \\
    \bottomrule
  \end{tabular}
\end{table}

Table~\ref{tab:ch4-heating-protection} shows that degraded and higher-current
conditions tripped earlier than low-current or healthy-socket conditions. The
result supports the protection objective because the relay response followed
the compensated estimated socket temperature, not merely current magnitude. A
limitation remains: the external NTC and compensation model still
underestimated the hottest terminal temperature in severe degraded-contact
cases. The prototype therefore demonstrates useful complementary protection,
but a final certified design would need a validated thermal model across more
socket types, contact materials, enclosure conditions, and ambient conditions.

